\documentclass{article}
\usepackage{svg}
\usepackage{amsmath}
\usepackage{geometry} % Required for customizing margins
\geometry{margin=1in} % Sets 1-inch margins on all sides

\title{ELCT 331 Control Systems Assignment 2}
\author{Jacob C. Vaught}
\date{Fall 2023}
\begin{document}
\maketitle

\section*{Problem \#1[20 Points]}
Reduce the block diagram shown in the figure to unity feedback form and find the system characteristic equation. 
\newline
\begin{figure}[htbp]
  \centering
  \includesvg[width=\textwidth]{331_HW2_P1.svg}
\end{figure}


\subsection*{Solution}
\[
\frac{Y(s)}{X(s)} = \frac{s+1}{s^3+3s^2+2s+1}
\]
The characteristic equation is given as:

\[
s^3 + 3s^2 + 2s + 1 = 0
\]
\newline
\begin{figure}[htbp]
  \centering
  \includesvg[width=\textwidth]{331_HW2_P1a.svg}
\end{figure}

\newpage
\section*{Problem \#2[20 Points]}
The Figure shows the block diagram of a dc-motor control system. The signal \( N(s) \) denotes torque at the motor shaft.

\begin{enumerate}
    \item[a. ] Find the transfer function \( H(s) \) so that the output \( Y(s) \) is not affected by the disturbance torque \( N(s) \).
    \item[b. ] With \( H(s) \) as determined in part (a), find the value of \( K \) so that the steady-state value of \( e(t) \) is equal to 0.1 when the input is a unit-ramp function \( r(t)=tu_s(t) \), \( R(s)=\frac{1}{s^2} \), and \( N(s)=0 \). Apply the final-value theorem.
\end{enumerate}
\newline
\begin{figure}[htbp]
  \centering
  \includesvg[width=\textwidth]{331_HW2_P2.svg}
  \Large \( G(s) = \frac{K(s + 3)}{s(s + 1)(s + 2)} \)
\end{figure}
\subsection*{Solutions}

\subsubsection*{(a) Deriving the Transfer Function}

Consider the given expression for \( Y(s) \):

\[
Y(s) = R(s) \cdot \frac{Y(s)}{R(s)} \bigg|_{N(s) = 0} + N(s) \cdot \frac{Y(s)}{N(s)} \bigg|_{R(s) = 0}
\]
It is stated that the transfer function \( \frac{Y(s)}{N(s)} \) when \( R(s) = 0 \) is:

\[
\frac{Y(s)}{N(s)} \bigg|_{R(s) = 0} = \frac{1}{1 + \frac{G(s)}{1 + G(s)H(s)}} = \frac{1 + G(s)H(s)}{1 + G(s)H(s) + G(s)}
\]
As per the given condition that \( Y(s) \) is not affected by \( N(s) \), we have:

\[
\frac{Y(s)}{N(s)} \bigg|_{R(s) = 0} = 0
\]
From this condition, it follows that \( 1 + G(s)H(s) = 0 \). Therefore, we can derive the expression for \( H(s) \):

\[
H(s) = -\frac{1}{G(s)} = -\frac{s(s+1)(s+2)}{K(s+3)}
\]

\subsubsection*{(b) Calculating the Gain \( K \)}
Given equations (1) and (2), when \( N(s) = 0 \), we can express \( E(s) \) as:

\[
E(s) = \frac{R(s)}{G(s)} = \frac{1/s^2}{\frac{K(s+3)}{s(s+1)(s+2)}} = \frac{(s+1)(s+2)}{Ks(s+3)}
\]
To find the steady-state error \( e(\infty) \), we take the limit:

\[
e(\infty) = \lim_{{s \to 0}} sE(s) = \lim_{{s \to 0}} \frac{(s+1)(s+2)}{K(s+3)} = \frac{2}{3K} = 0.1
\]
Solving for \( K \), we get:

\[
K = \frac{20}{3} = 6.67
\]


\newpage
\section*{Problem \#3[20 Points]}
\begin{enumerate}
    \item[a. ] Draw an equivalent Signal Flow Graph (SFG) for the system.
    \item[b. ] Find the following transfer functions by applying the gain formula of the SFG directly to the block diagram.
    \[
    \frac{Y(s)}{R(s)} \bigg|_{N=0} \quad \frac{Y(s)}{N(s)} \bigg|_{R=0} \quad \frac{E(s)}{R(s)} \bigg|_{N=0} \quad \frac{E(s)}{N(s)} \bigg|_{R=0}
    \]
    \item[c. ] Compare the answers by applying the gain formula to the equivalent SFG.
\end{enumerate}
\newline
\begin{figure}[htbp]
  \centering
  \includesvg[width=\textwidth]{331_HW2_P3.svg}
\end{figure}

\subsection*{Solution}
\subsubsection*{(a) Drawing the Signal Flow Graph}
\newline
\begin{figure}[htbp]
  \centering
  \includesvg[width=0.5\textwidth]{331_HW2_P3a.svg}
\end{figure}
\newpage
\subsubsection*{(b) Transfer Functions}

We consider two feedback loops: one with a gain of \(-G_3H_1\) and another with a gain of \(-G_2G_3\). Since there are no non-touching loops, we can calculate the determinant \(\Delta\) as:

\[
\Delta = 1 + G_3H_1 + G_2G_3
\]

\paragraph{Expression for \( \frac{Y(s)}{R(s)} \) under \( N(s) = 0 \)}

Here, we have two individual loop gains \(M_1 = G_2G_3\) and \(M_2 = G_1G_3\), and their corresponding \(\Delta_k\) values are both 1. The overall transfer function from \( R(s) \) to \( Y(s) \) under \( N(s) = 0 \) can be written as:

\[
\frac{Y(s)}{R(s)} \bigg|_{N(s) = 0} = \frac{1}{\Delta} \sum_k M_k \Delta_k = \frac{G_2(s)G_3(s) + G_1(s)G_3(s)}{1 + G_3(s)H_1(s) + G_2(s)G_3(s)}
\]

\paragraph{Expression for \( \frac{Y(s)}{N(s)} \) under \( R(s) = 0 \)}

Here, we only have a single loop with a gain of \(M_1 = 1\) and \(\Delta_1 = 1\). Thus, the transfer function \( \frac{Y(s)}{N(s)} \) when \( R(s) = 0 \) is:

\[
\frac{Y(s)}{N(s)} \bigg|_{R(s) = 0} = \frac{1}{\Delta} \sum_k M_k \Delta_k = \frac{1}{1 + G_3(s)H_1(s) + G_2(s)G_3(s)}
\]

\paragraph{Expression for \( \frac{E(s)}{R(s)} \) under \( N(s) = 0 \)}

In this scenario, \( M_1 = 1 \) with \(\Delta_1 = 1 + G_3H_1\) and \( M_2 = -G_1G_3 \) with \(\Delta_2 = 1\). Therefore, the transfer function \( \frac{E(s)}{R(s)} \) when \( N(s) = 0 \) is:

\[
\frac{E(s)}{R(s)} \bigg|_{N(s) = 0} = \frac{1}{\Delta} \sum_k M_k \Delta_k = \frac{1 + G_3(s)H_1(s) - G_1(s)G_3(s)}{1 + G_3(s)H_1(s) + G_1(s)G_3(s)}
\]

\paragraph{Expression for \( \frac{E(s)}{N(s)} \) under \( R(s) = 0 \)}

Here, \( M_1 = -1 \) with \(\Delta_1 = 1\). Thus, the transfer function \( \frac{E(s)}{N(s)} \) under \( R(s) = 0 \) is:

\[
\frac{E(s)}{N(s)} \bigg|_{R(s) = 0} = \frac{1}{\Delta} \sum_k M_k \Delta_k = \frac{-1}{1 + G_3(s)H_1(s) + G_2(s)G_3(s)}
\]


\newpage
\section*{Problem \#4 [30 points]}

The figure shows the block diagram of a control system with conditional feedback. The transfer function \( G_p(s) \) denotes the controlled process, and \( G_c(s) \) and \( H(s) \) are the controller transfer functions.

\begin{enumerate}
    \item[a. ] Derive the transfer functions 
    \[
    \frac{Y(s)}{R(s)} \bigg|_{N=0}  \quad \frac{Y(s)}{N(s)} \bigg|_{R=0}
    \]
    Find \( \frac{Y(s)}{R(s)} \bigg|_{N=0} \) when \( G_c(s) = G_p(s) \).
    
    \item[b. ] Let
    \[
    G_p(s) = G_c(s) = \frac{100}{(s+1)(s+5)}
    \]
    Find the output response \( y(t) \) when \( N(s) = 0 \) and \( r(t) = u_s(t) \).

    \item[c. ] With \( G_p(s) \) and \( G_c(s) \) as given in part (b), select \( H(s) \) among the following choices such that when \( n(t) = u_s(t) \) and \( r(t) = 0 \), the steady-state of \( y(t) \) is equal to zero. (There may be more than one answer.)
\begin{center}
  \begin{minipage}{0.4\textwidth}
    \begin{align*}
      H(s) &= \frac{10}{s(s+1)} \\
      H(s) &= \frac{10(s+1)}{s+2}
    \end{align*}
  \end{minipage}
  \begin{minipage}{0.4\textwidth}
    \begin{align*}
      H(s) &= \frac{10}{(s+1)(s+2)} \\
      H(s) &= \frac{K}{s^n} \text{ (n is a positive integer. Select n.)}
    \end{align*}
  \end{minipage}
\end{center}

    
    
    Keep in mind that the poles of the closed-loop transfer function must all be in the left-half \( s \)-plane for the final-value theorem to be valid.
\end{enumerate}
\newline
\begin{figure}[htbp]
  \centering
  \includesvg[width=\textwidth]{331_HW2_P4.svg}
\end{figure}
\subsection*{Solution}

\subsubsection*{(a) Transfer Functions for Various Cases}

First, let's define \( \delta = 1 + G_p H \).

\paragraph{For \( \frac{Y(s)}{R(s)} \) under \( N(s) = 0 \)}
For this case, \( M_1 = G_p \) and \( \delta_1 = 1 \), and \( M_2 = G_c H G_p \) and \( \delta_2 = 1 \).
The transfer function is:
\[
\frac{Y(s)}{R(s)} \bigg|_{N(s) = 0} = \frac{1}{\delta} \sum_k M_k \delta_k = \frac{G_p(s) + G_c(s) H(s) G_p(s)}{1 + G_p(s) H(s)}
\]

\paragraph{For \( \frac{Y(s)}{N(s)} \) under \( N(s) = 0 \)}
In this scenario, \( M_1 = G_p \) and \( \delta_1 = 1 \).
The transfer function becomes:
\[
\frac{Y(s)}{N(s)} \bigg|_{N(s) = 0} = \frac{1}{\delta} \sum_k M_k \delta_k = \frac{G_p(s)}{1 + G_p(s) H(s)}
\]

\paragraph{For \( \frac{Y(s)}{R(s)} \) under \( N(s) = 0 \) when \( G_c(s) = G_p(s) \)}
Here, \( M_1 = G_p \) and \( \delta_1 = 1 \), and \( M_2 = H G_p^2 \) and \( \delta_2 = 1 \).
The transfer function is:
\[
\frac{Y(s)}{R(s)} \bigg|_{N(s) = 0} = \frac{1}{\delta} \sum_k M_k \delta_k = \frac{G_p(s)[1 + H(s) G_p(s)]}{1 + G_p(s) H(s)}
\]

\subsubsection*{(b) Calculations and Partial Fraction Decomposition}

For \( Y(s) \) under \( N(s) = 0 \), we have:
\[
Y(s) \bigg|_{N(s) = 0} = \frac{Y(s)}{R(s)} \bigg|_{N(s) = 0} R(s) = \frac{G_p(s) R(s)}{1 + G_p(s) H(s)} = \frac{100}{s(s+1)(s+5)}
\]
Performing partial fraction decomposition, the result is:
\[
Y(s) = \frac{20}{s} - \frac{25}{s+1} + \frac{5}{s+5}
\]
Taking the inverse Laplace transform, we get:
\[
y(t) = 20 - 25e^{-t} + 5e^{-5t}, \quad \text{for } t \geq 0
\]

\subsubsection*{(c) Selecting \( H(s) \) for Zero Steady-State Output}

To determine which \( H(s) \) will make the steady-state output \( y(t) \) zero for \( n(t) = u_s(t) \) and \( r(t) = 0 \), we evaluate \[ \lim_{{s \to 0}} G_p(s)H(s) \] and check if it equals -1. From part (b) that \( G_p(s) = \frac{100}{s(s+1)(s+5)} \).

We examine each of the \( H(s) \) options:

\begin{enumerate}
    \item \( H_1(s) = \frac{10}{s(s+1)} \)
    \[ \lim_{{s \to 0}} \frac{100}{s(s+1)(s+5)} \times \frac{10}{s(s+1)} = \infty \neq -1 \]
    Hence, \( H_1(s) \) does not satisfy the condition.

    \item \( H_2(s) = \frac{10(s+1)}{s+2} \)
    \[ \lim_{{s \to 0}} \frac{100}{s(s+1)(s+5)} \times \frac{10(s+1)}{s+2} = 0 \neq -1 \]
    Hence, \( H_2(s) \) does not satisfy the condition.

    \item \( H_3(s) = \frac{10}{(s+1)(s+2)} \)
    \[ \lim_{{s \to 0}} \frac{100}{s(s+1)(s+5)} \times \frac{10}{(s+1)(s+2)} = \infty \neq -1 \]
    Hence, \( H_3(s) \) does not satisfy the condition.

    \item \( H_4(s) = \frac{K}{s^n} \)
    \[ \lim_{{s \to 0}} \frac{100}{s(s+1)(s+5)} \times \frac{K}{s^n} = \lim_{{s \to 0}} \frac{100K}{s^{n+1}(s+1)(s+5)} \]
    To make this limit -1, \( n+1 \) must equal 1 so that the \( s^{n+1} \) term cancels out, resulting in a finite limit. The choice of \( K \) would then be such that the limit is -1.
\end{enumerate}

Conclusively, none of the given \( H(s) \) options would make \( y(t) \) approach zero in the steady-state, given the \( G_p(s) \) specified in part (b).


\end{document}
